\section{A posterior certificate using robustness only where it is needed}
\label{sec:pointwise}

The earlier compression proof uses robustness only at its final prediction
comparison. This permits a version that does not assume robust support.
Fix a radius $0<r\le1$. Define
\[
 b_w(x;r)=\mathbf1\{\text{some perturbation of total RMS norm at most }r
 \text{ makes }\sgn z_w(x)\,z_w(x;\eta)\le0\}.
\]
Clean zeros are declared bad. For a parameter law $Q$, write
$b_Q(x;r)=\E_Q b_w(x;r)$, and let $b_D(Q;r)$ and $b_s(Q;r)$ be its
population and empirical input averages. The reference measure $D$ here is
the input distribution, not a target-mass exponent.

\begin{proposition}[Compression with a pointwise defect]
For $D_{\rm KL}(Q\Vert\mu_{N,t})\le NK/t^2$, the same finite class
$\mathcal H(d,L,K,r,\alpha)$ from the earlier robust entropy-compression
theorem admits a probability law $\Gamma_Q$ with
\[
 \left|\E_Q\mathbf1\{\sgn z_w(x)\ne f(x)\}
       -\E_{h\sim\Gamma_Q}\mathbf1\{h(x)\ne f(x)\}\right|
 \le\alpha+b_Q(x;r)\qquad(x\in\{-1,1\}^d).
\]
The construction and the class do not depend on the target $f$.
\end{proposition}
\begin{proof}
Run the earlier construction for the unrestricted law $Q$. Its entropy
budgets, clipping bounds, moment estimates, and common candidate marginal
do not use robustness. At each fixed input, its comparison coupling has
the original full-network marginal $Q$ and, except with probability less
than $\alpha/4$, produces an exact postactivation perturbation of norm
less than $r/2$. When $b_w(x;r)=0$, the candidate has the original sign
and margin greater than $r/2$: otherwise an additional output injection
within the remaining norm budget could reach zero. The same finite-grid
rounding preserves that sign. In the small-width branch, good parameter
clipping likewise preserves the original network until rounding; local
robustness gives its output margin greater than $r$. Adding the probability
$b_Q(x;r)$ to the original exceptional probability proves the assertion.
No conditioning on a robustness event is performed, so no extra entropy
cost or inconsistency between pointwise couplings is introduced.
\end{proof}

\begin{theorem}[Robustness-defect certificate for a consistent posterior]
With probability at least $1-\delta$ over $n$ iid labeled examples,
simultaneously for all $N,t,Q$ with the stated entropy bound,
\[
 \operatorname{GibbsErr}_D(Q)
 \le 2\operatorname{GibbsErr}_s(Q)+3\alpha+b_D(Q;r)+2b_s(Q;r)
       +\frac{8[\log|\mathcal H|+\log(1/\delta)]}{n}.
\]
In particular the empirical Gibbs error term vanishes for the original
noiseless-consistent Gaussian posterior. This statement imposes no
robustness restriction on that prior.
\end{theorem}
\begin{proof}
On the usual finite-class relative-Chernoff event,
$\operatorname{err}_D(h)\le2\operatorname{err}_s(h)
+8[\log|\mathcal H|+\log(1/\delta)]/n$ for every $h\in\mathcal H$.
Average under the common law $\Gamma_Q$ and use the preceding proposition
once on the population and once on the empirical inputs.
\end{proof}

For a fixed posterior selected without using an independent unlabeled
sample $U$ of size $m$, Hoeffding also gives
\[
 b_D(Q;r)\le b_U(Q;r)+\sqrt{\log(1/\delta')/(2m)}
\]
with probability at least $1-\delta'$. Sound upper certificates for
$b_w(x;r)$ may replace its exact value. This is a mathematical posterior
certificate; one trained network is not automatically a draw from the
specified Bayesian posterior. The training defect cannot simply be omitted:
the posterior depends on the training inputs, so small population defect
alone does not establish a small empirical defect.

\subsection{A distributional prior with robust training likelihood}
There is also a genuine learning theorem under a weaker prior restriction.
Suppose the input distribution $D$ is known before the labels, set
$J=L-1$, $E_v=N\|v\|^2$, and put
\[
 r_{\rm ref}(w)=\frac{1}{100\sqrt{1+J E_v}}.
\]
For fixed $0\le\rho<1$, condition the raw Gaussian only on
\[
 D\{x:b_w(x;r_{\rm ref}(w))=1\}\le\rho.
\]
Unlike the earlier prior, this condition imposes no all-input clean margin.
For every observed example $(x_i,y_i)$, use the likelihood requiring
$y_i z_w(x_i;\eta)>0$ for every perturbation of norm at most
$r_{\rm ref}(w)$. Thus the likelihood enforces a robust answer on training
inputs. This is a modified learning rule, not the original sign-only
likelihood with a silently omitted empirical defect.

\begin{corollary}[Polynomial-depth learning up to a distributional defect]
For targets with a width-$C$, depth-$(L-1)$ layered Boolean circuit, this
posterior has Gibbs error at most $\rho+\varepsilon$ after
$\operatorname{poly}(d,L,C,\varepsilon^{-1},\log(1/\delta))$ examples,
simultaneously for all $N\ge C$ and all $0<t\le1$. The polynomial has
degree independent of depth, width, and Gaussian scale.
\end{corollary}
\begin{proof}
The earlier robust target cylinder belongs to this prior support and
satisfies the robust training likelihood for every sample. If its raw
mass is at least $\exp(-ND_*/t^2)$, posterior KL is at most $ND_*/t^2$.
With $B=\max\{8,8[D_*+\log(1/\beta)]\}$, the same Gaussian readout tail
bound gives posterior mass at most $\beta$ to $E_v>B$. On the complementary
event the required radius is at least
$r_0=1/[100\sqrt{1+JB}]$. The conditional posterior has entropy bound
$N(D_*+1)/t^2$, population robustness defect at most $\rho$, and empirical
defect zero. Apply the preceding theorem and add the discarded mass:
\[
 \operatorname{GibbsErr}_D(Q_s)
 \le\beta+3\alpha+\rho
       +8[\log|\mathcal H|+\log(1/\delta)]/n.
\]
Choosing $\beta=\varepsilon/4$, $\alpha=\varepsilon/12$ and the earlier
polynomial sample size proves the claim. The finite class is independent
of $N,t$, which gives one simultaneous sample event.
\end{proof}

If a fixed reference input distribution $\nu$ is used instead, and
$D\le M\nu$, replace the population defect $\rho$ by $M\rho$. This makes
explicit what distributional information enters the prior. The result
does not assert that the support can be estimated from arbitrary unlabeled
data without an additional uniform-convergence argument.
